\section{Data Characteristics and Challenges}

\subsection{Forest Masking Using ESA WorldCover}
Because NDVI measures relative greenness but does not distinguish among vegetation types, isolating forested pixels is essential when evaluating forest health. The \textit{ESA WorldCover} 10\,m global land cover product provides a high-resolution classification that we use to identify and mask non-forest areas prior to vegetation index computation and time-series analysis. Without such land-cover context, signals from grasslands, croplands, water bodies, or built environments could be misinterpreted as forest dynamics, confounding anomaly detection. This practice of restricting analysis to true forest pixels is also employed by Ghilardi et al.\ \cite{worldcover} in their multi-decadal forest disturbance study, where forest masking improves the accuracy of canopy loss and recovery metrics derived from NDVI. By incorporating WorldCover early in the processing pipeline, we reduce contamination from non-forest surfaces, improve signal interpretability, and focus computational resources on pixels that are ecologically relevant to our disease anomaly analysis.